\section{An autonomous processor with a spatial clock}
\label{sec:sa-autonomous}
The preceding constructions repeat a time-dependent pulse.  We now make
that repetition part of an autonomous divergence-free field: a third
coordinate runs through the phases of a planar Hamiltonian motion.
The observer is restricted to a quiet clock window, so transient transport
paths cannot give false detections.  An onto scalar time change then starts
the fluid from rest and makes its force decay.

\begin{theorem}\label{thm:sa-autonomous}
At fixed positive computable viscosity, a deterministic one-tape machine
and finite input effectively determine a smooth mean-zero general force
on the unit torus such that, for every $k\ge0$ and spatial multi-index
$\beta$,
\begin{equation}\label{eq:sa-autonomous-force-class}
 \|\partial_t^k\partial_x^\beta f(t)\|_\infty
       \le C_{k,\beta}(1+t)^{-1}.
\end{equation}
Its zero-data Navier--Stokes solution is globally smooth and unique in
the classical periodic comparison class, with uniformly bounded kinetic
energy and mean-zero pressure.  For the fixed label and open set
\begin{equation}\label{eq:sa-autonomous-event}
 y_*=(1/2,1/2,1/4),\qquad
 \mathcal O=\{(x_1,x_2,z):1/2<x_1<1,\ 0<z<1/8\},
\end{equation}
its material trajectory enters $\mathcal O$ exactly when the machine
halts.  Samples encode the full machine computation through the first
halt, or indefinitely if there is none.  The prescription uses only the
finite table and input; arbitrary fixed real positive viscosity is allowed
as a relative coefficient.
\end{theorem}

\subsection{A recorder with identifiable incoming instructions}
\label{subsec:sa-autonomous-processor}
Normalize to one halt state $q_f$ and a total nonhalt table.  Add a fresh
initial state $q_I$ whose instructions preserve the symbol, stay put,
and enter the old initial state.  No rule enters $q_I$.  This guarantees
a positive number of steps before any halt and ensures that the initial
control is never visited again.  Let $Q$ and $A$ be the normalized
state set and work alphabet.
Apply the frontier recorder of Lemma~\ref{lem:sb-frontier} to this
normalized ordinary machine, using exactly its notation:
\[
 \mathcal R=(Q\setminus\{q_f\})\times A,\qquad
 \Sigma=A\times\{\mathtt e,\mathtt F,\lambda_r:r\in\mathcal R\}
                  \times\{0,1\}.
\]
The main controls are $B_q$; the other controls are $A_r,R_r,N_q,L_q$.
At a main checkpoint the work track represents the ordinary tape, all
return marks vanish, and a rightward frontier follows the finite history
block. One ordinary step writes and moves the work head, marks its new
position, records the instruction at the frontier, advances that frontier,
and returns to the marked position. Lemma~\ref{lem:sb-frontier} proves
this cycle and its full-domain inverse with every nonfrontier guard.
In particular each target control has a fixed incoming displacement and
the written whole symbol determines its incoming instruction. These are
the precise hypotheses used for the reciprocal rectangles below.

Initialize at $B_{q_I}$, head zero, with the input work tape, all marks
zero, history frontier $\mathtt F$ at cell two and $\mathtt e$ elsewhere.
At the $n$th ordinary checkpoint, after primitive count $m_n$, the work
configuration is the normalized machine's $n$th configuration, records
occupy cells $2,\ldots,n+1$, and the frontier is at $n+2$.
If $h_{n+1}$ is the new work-head position, the same lemma gives
\begin{equation}\label{eq:sa-autonomous-count}
       m_{n+1}-m_n=2(n+2-h_{n+1})+4.
\end{equation}
Thus each ordinary step takes finitely many positive primitive steps.
The fresh state $q_I$ occurs only at the initial ordinary checkpoint,
and at least one ordinary step precedes any halt. In a nonhalting run
no positive primitive sample has main control $B_{q_f}$ or $B_{q_I}$.

\subsection{Reciprocal rectangle maps}
\label{subsec:sa-autonomous-rectangles}
Use odd digits and base $K=2|\Sigma|+1$ in Lemma~\ref{lem:sa-radix}.
For a tape relative to its head let
$x=E(\sigma_{-1}\sigma_{-2}\cdots)$ and
$y=E(\sigma_0\sigma_1\cdots)$, and encode by $o_r+(x,y)$.
The two internal initial coordinates $x_{\rm in},y_{\rm in}$ are
rational.  Choose
\begin{equation}\label{eq:sa-autonomous-offsets}
 o_{B_{q_I}}=(-x_{\rm in},-y_{\rm in}),\qquad
 o_{B_{q_f}}=(2,0),\qquad o_r=(-3j,0)\quad(j=1,2,\ldots)
\end{equation}
for the remaining controls.  The initial code is zero, and the closed
state squares $o_r+[0,1]^2$ are disjoint.

For an instruction $(r,\eta)\mapsto(r',\beta,d)$ write
$h_\sigma(v)=(d(\sigma)+v)/K$.  The normalized branches are
\begin{equation}\label{eq:sa-autonomous-branches}
\begin{array}{ll}
 d=0:&(x,h_\eta(v))\mapsto(x,h_\beta(v)),\\
 d=1:&(x,h_\eta(v))\mapsto(h_\beta(x),v),\\
 d=-1:&(h_\gamma(w),h_\eta(v))\mapsto
          (w,h_\gamma(h_\beta(v))),\quad\gamma\in\Sigma.
\end{array}
\end{equation}
All free variables range independently over $[0,1]$.
These are exactly the head-relative list operations.  Sources are
separated by control, read symbol, and the extra left symbol for a left
move.  For a common target control the incoming move is fixed and distinct
rules have distinct written symbols.  Those symbols separate the right
first digit for a stay, the left first digit for a right move, and the
right pair $(\gamma,\beta)$ for a left move.  Thus targets are separately
separated as well.  The linear parts are $I$, $\diag(1/K,K)$, and
$\diag(K,1/K)$, respectively.  Unlike general prefix maps, all preserve
planar area.  This leaves the third coordinate available for a clock.

\subsection{A Hamiltonian extension without planar contraction}
\label{subsec:sa-autonomous-extension}
The area-preserving requirement rules out the simultaneous contraction
used in Lemma~\ref{lem:sa-planar-routing}.  Instead we separate the
rectangle centers, carry the unchanged rectangles to distant parking
positions, and deform them there with reciprocal scales.

\begin{lemma}\label{lem:sa-autonomous-extension}
Let $U_i,\widehat U_i$, $1\le i\le J$, $J\ge1$, be two separately
disjoint families of nondegenerate rational closed rectangles in
$\mathbb R^2$.  For each pair prescribe
\[
 X\mapsto\widehat c_i+
       \diag(\lambda_i,\lambda_i^{-1})(X-c_i),\qquad\lambda_i>0
\]
with rational $\lambda_i$ and the respective centers, and require the
compatibility condition
\[
 \widehat U_i=\widehat c_i+
       \diag(\lambda_i,\lambda_i^{-1})(U_i-c_i).
\]
There is an effective smooth compactly supported planar Hamiltonian
field $V(\xi,X)$, with time support in $(1/4,3/4)$, whose unit-time
map has this formula on each $U_i$ and a positive neighborhood.
The construction gives a rational bound on its spatial support and all
controlled rectangle paths.
\end{lemma}
\begin{proof}
Choose a positive rational margin $\epsilon$ at most one quarter of a
positive coordinate-separating gap for each pair in each family; for
empty pair lists take $\epsilon=1$.  The enlarged families stay disjoint.
Choose rational $D_0$ larger than $\epsilon$ plus every source or target
half-side.  Take an integer $G\ge1$ so that within each family the
scaled centers $Gc_i$ and $G\widehat c_i$ are more than $10D_0$ apart
in the sup norm.  Choose parking centers $e_i$ more than $10D_0$ apart
from one another and from all these scaled centers.  A sufficiently
distant row with spacing $20D_0$ does this effectively.

Use affine paths
\begin{equation}\label{eq:sa-autonomous-moving-box}
 X\mapsto C_i(\xi)+\diag(a_i(\xi),a_i(\xi)^{-1})(X-c_i).
\end{equation}
First translate centers simultaneously $c_i\to Gc_i$, keeping shapes
fixed; any coordinate separation can only increase.  Next move one
rectangle at a time to $e_i$.  At those parking centers change $a_i$
from $1$ to $\lambda_i$.  Each half-side stays between its endpoint
values and hence below $D_0-\epsilon$.  Move the target-shaped rectangles
one at a time to $G\widehat c_i$, then translate centers simultaneously
down to $\widehat c_i$.  In this last phase the target separations stay
valid down to scale one.
During source evacuation, occupied centers are scaled source centers or
parking centers. During target delivery, every source has already been
evacuated, and occupied centers are parking centers or scaled target
centers. Each such mixed family has separation greater than $10D_0$ by
the choices above. Thus arbitrary intersections between the original
source and target families never place two occupied rectangles at the
same location. The obstacle family for each transfer consists of every
other currently occupied rectangle, wherever it is in that phase.

For a one-at-a-time transfer, put a closed square of sup radius $4D_0$
around every stationary center.  These obstacles are disjoint, and the
transfer endpoints are outside them.  Follow the straight segment, replacing
any portion inside one obstacle by a boundary detour in a fixed orientation.
Entry and exit points solve rational linear inequalities, so all waypoints
are rational.  A tangent contact, a corner contact without entry, or a segment already
on the boundary needs no detour: it already has the required clearance.
A segment crossing the interior has rational first and last intersection
points, found by clipping its rational parameter interval against the
four rational side inequalities. A fixed clockwise boundary arc between
those points uses only rational corners. The obstacles are separated,
so each detour misses every other obstacle.  The moving center remains at sup distance at least
$4D_0$ from all stationary centers; both enlarged rectangles have radius
less than $D_0$, proving separation throughout.

Schedule the finitely many submotions in successive equal slots inside
$[1/3,2/3]$, using a flat step for each segment and scale change.  This
smooths time without changing any path or clearance.  Let $\chi_i$
equal one on a positive collar of the moving rectangle and be supported
in its $\epsilon$ enlargement.  With $Y=X-C_i$, set
\begin{equation}\label{eq:sa-autonomous-hamiltonian}
 \mathcal H=\sum_i\chi_i\left(
       \dot C_{i1}Y_2-\dot C_{i2}Y_1
                  +\frac{\dot a_i}{a_i}Y_1Y_2\right),\qquad
 V=(\partial_{X_2}\mathcal H,-\partial_{X_1}\mathcal H).
\end{equation}
The supports are disjoint.  On rectangle $i$ the field is
$\dot C_i+\diag(\dot a_i/a_i,-\dot a_i/a_i)(X-C_i)$, so
\eqref{eq:sa-autonomous-moving-box} solves its ODE exactly.
The same formula holds on a positive initial neighborhood by the plateau
and finite affine norm bounds.  Smooth ODE uniqueness identifies the
motion.  A rational support bound follows from the finite endpoint and
waypoint list plus $D_0+1$.  Every operation is effective finite geometry
and the fixed flat step, proving the lemma.
\end{proof}

Apply this lemma to \eqref{eq:sa-autonomous-branches}. With
$I_\eta=h_\eta([0,1])$, the normalized target rectangles for moves
$0,1,-1$ are, respectively, $[0,1]\times I_\beta$,
$I_\beta\times[0,1]$, and $[0,1]\times h_\gamma(I_\beta)$.
Their side lengths are exactly those of their source rectangles multiplied
by $(1,1)$, $(K^{-1},K)$, and $(K,K^{-1})$. Adding the respective state
offsets verifies the compatibility equation for every branch.
The branch list is
nonempty because of the fresh initial state.  Choose rational $M>1$ so
that its support, controlled paths, and all state squares lie in $(-M,M)^2$.
We now have one complete planar transition pulse with its own proof of
area preservation; no third-dimensional strain has been used.

\subsection{The spatial clock and the observation window}
\label{subsec:sa-autonomous-clock}
Set $\kappa=1/(16M)$ and $T(X)=(1/2,1/2)+\kappa X$.
All relevant geometry lies after scaling in $(7/16,9/16)^2$.
Periodically extend
\[
 \widetilde V(x_1,x_2,z)=\kappa V(z,T^{-1}(x_1,x_2)).
\]
Its known zero collars at the horizontal boundaries and at $z=0,1$
make this smooth.  It has horizontal divergence zero; each horizontal
component has integral zero in $(x_1,x_2)$ since it is a scaled derivative
of the compact Hamiltonian.

For an explicit vertical clock choose a flat step $\theta$ that is zero
on $(-\infty,1/3]$ and one on $[2/3,\infty)$.  Put
\[
 \psi(s)=(1-\theta(16(s-1/8)))
             (1-\theta(16(-s-1/8))),\qquad
 b(x_1,x_2)=\psi(x_1-1/2)\psi(x_2-1/2),
\]
periodize $b$, and set $w(x_1,x_2)=b(x_1,x_2)-b(x_1-1/2,x_2)$.
The bump is one on $[3/8,5/8]^2$ and supported in
$(5/16,11/16)^2$.  Its half-period translate has disjoint support,
so $w=1$ on all scaled state squares and controlled paths, and $\int w=0$.
The autonomous field
\begin{equation}\label{eq:sa-autonomous-W}
 W=(\widetilde V_1,\widetilde V_2,w)
\end{equation}
is divergence free and mean zero.  Write $\Psi_\tau$ for its global
smooth flow on the torus.

From $(T(X),1/4)$ with $X$ in a source rectangle, the clock increases
at speed one and the horizontal coordinates follow its prescribed path
with planar phase $\xi=1/4+\tau$.  After that pulse, the horizontal target
stays fixed while the clock passes from $z=0$ to $z=1/4$.  The entire
path has $w=1$, so this description solves the autonomous ODE.  Thus at
unit times the map performs one primitive transition and returns the clock
to $1/4$.  If $X_j$ is the primitive configuration code, induction gives
\begin{equation}\label{eq:sa-autonomous-samples}
       \Psi_j(y_*)=(T(X_j),1/4)
\end{equation}
through halt or indefinitely in a nonhalting run.
Along these controlled paths, $z=1/4+\tau\pmod1$.  Hence for every
positive such $j$, in the window
\begin{equation}\label{eq:sa-autonomous-window}
                  j-1/4<\tau<j-1/8
\end{equation}
the horizontal coordinate is exactly $T(X_j)$ and $0<z<1/8$.
If the machine halts, its primitive arrival index is positive and its
halt square has first offset $2$, so this window meets $\mathcal O$.
If it never halts, the clock test in \eqref{eq:sa-autonomous-event} holds
only in these windows.  At their positive samples the control is neither
$B_{q_I}$ nor $B_{q_f}$, hence its square has first offset $-3k$, $k\ge1$,
and $T(X_j)_1<1/2$.  This excludes all false detections, regardless of
where intermediate transport paths pass.
The samples at $m_n$ encode the normalized work configuration.
Discarding its initial stay step recovers the supplied machine's
computation.  Count the records to recover $n$, including $n=0$ at
initialization.  The frontier has absolute position $n+2$; subtracting
its decoded head-relative index recovers the absolute head position.
The work track then gives the absolute tape configuration.

\subsection{Starting from rest while covering all internal time}
\label{subsec:sa-autonomous-force}
Use the scalar profile
\begin{equation}\label{eq:sa-autonomous-profile}
 A_0(t)=\frac{t}{(1+t)^2}=(1+t)^{-1}-(1+t)^{-2}
\end{equation}
and prescribe
\begin{equation}\label{eq:sa-autonomous-force}
 f=A_0'W+A_0^2(W\cdot\nabla)W-\nu A_0\Delta W.
\end{equation}
Then $(u,p)=(A_0W,0)$ solves the forced equation and has $u(0)=0$.
In contrast, $A_0'(0)=1$ gives $f(0)=W$; the force need not vanish at
zero.  All three terms have zero spatial mean, using
$(W\cdot\nabla)W=\operatorname{div}(W\otimes W)$.
Lemma~\ref{lem:p2-realization} gives uniqueness in the periodic comparison
class.  Since $0\le A_0\le1/4$, the kinetic energy is at most
$\|W\|_2^2/32$.

For explicit derivative bounds put
$B_{r,k}=r(r+1)\cdots(r+k-1)$, with $B_{r,0}=1$.
The $k$th derivative of $(1+t)^{-r}$ has absolute value
$B_{r,k}(1+t)^{-r-k}$.
Writing
\[
 a_k=B_{1,k}+B_{2,k},\quad b_k=B_{2,k}+2B_{3,k},\quad
 c_k=B_{2,k}+2B_{3,k}+B_{4,k},
\]
and letting $M_{0,\beta},M_{1,\beta},M_{2,\beta}$ be the sup norms of
the spatial $\beta$ derivatives of $W,(W\cdot\nabla)W,\Delta W$,
respectively, gives
\[
 \|\partial_t^k\partial_x^\beta f(t)\|_\infty
 \le (b_kM_{0,\beta}+c_kM_{1,\beta})(1+t)^{-2-k}
              +\nu a_kM_{2,\beta}(1+t)^{-1-k}.
\]
This proves \eqref{eq:sa-autonomous-force-class} at every order, including
time zero, and supplies the sharper derivativewise decay as well.

The material flow is
\begin{equation}\label{eq:sa-autonomous-material}
 F(t,y)=\Psi_{S(t)}(y),\qquad
 S(t)=\int_0^t A_0(s)\,ds=\log(1+t)+(1+t)^{-1}-1.
\end{equation}
Differentiation proves the formula.  Since $A_0(t)>0$ for $t>0$ and
$S(t)\to\infty$, this clock is continuous and strictly increasing from
zero onto $[0,\infty)$.  Every finite autonomous event and every sample
$S^{-1}(j)$ therefore occur at finite physical time.  The established
window equivalence proves the theorem.

All finite tables, initial geometric tails, rectangle offsets, clearances,
waypoints, and scaling constants are effective.  The Hamiltonian cutoff
uses only rescaled flat steps and positive rational margins, and its
periodization has known zero collars.  The remaining force operations
are finite differentiation and the rational profile \eqref{eq:sa-autonomous-profile}.
The input affects the initial-state translation in
\eqref{eq:sa-autonomous-offsets}; no list of executed states is used.
At every finite time the material map is a smooth diffeomorphism, with
$u=\partial_tF\circ F^{-1}$ and
$\partial_t^2F=(\nu\Delta u-\nabla p+f)\circ F$, as follows by
 differentiating its trajectory equation.
